\subsection*{Problema 16}

\textbf{Enunciado:}
Calcular la integral indefinida:
\begin{equation*}
  I = \int \dfrac{dx}{1 + \sin x + \cos x}
\end{equation*}

\textbf{Solución:}
Para resolver esta integral trigonométrica, aplicamos la sustitución universal de Weierstrass, definida como:
\begin{equation*}
  z = \tan\left(\dfrac{x}{2}\right), \quad -\pi < x < \pi
\end{equation*}

A partir de esta sustitución, las identidades trigonométricas fundamentales en términos de $z$ se expresan como:
\begin{align*}
  \sin x &= \dfrac{2z}{1 + z^2}, \quad 
  \cos x = \dfrac{1 - z^2}{1 + z^2} \\
  dx &= \dfrac{2 \, dz}{1 + z^2}
\end{align*}

Sustituimos estas expresiones en la integral original $I$:
\begin{align*}
  I &= \int \dfrac{\dfrac{2 \, dz}{1 + z^2}}{1 + \dfrac{2z}{1 + z^2} + \dfrac{1 - z^2}{1 + z^2}}
\end{align*}

Simplificamos el denominador multiplicando toda la fracción por $(1 + z^2)$:
\begin{align*}
  I &= \int \dfrac{2 \, dz}{(1 + z^2) + 2z + (1 - z^2)} \\
  I &= \int \dfrac{2 \, dz}{2 + 2z} = \int \dfrac{2 \, dz}{2(1 + z)}
\end{align*}

Cancelamos el factor común $2$:
\begin{align*}
  I &= \int \dfrac{dz}{1 + z}
\end{align*}

La integración directa nos da el logaritmo natural:
\begin{align*}
  I &= \ln|1 + z| + C
\end{align*}

Finalmente, regresamos a la variable original $x$ sustituyendo $z = \tan(x/2)$:
\begin{equation*}
  \int \dfrac{dx}{1 + \sin x + \cos x} = \ln\left|1 + \tan\left(\dfrac{x}{2}\right)\right| + C
\end{equation*}

\begin{figure}[H]
	\centering
	\includegraphics[width=\columnwidth]{semana_04/graficas/SEM04_P16_grafica.png}
	\caption{Representación gráfica de $f(x)$ y su antiderivada $F(x)$ evaluada para la constante de integración $C = 0$ (Problema 16).}
	\label{fig:problema16}
\end{figure}
